\section{Conclusion}
\label{sec:conclusion}

We set out to evaluate a projective polynomial activation --- $f(z) = z +
\lambda z^{3}$ followed by an $\ell_2$ projection --- against the activation
functions and normalization layers in current use, and to separate the two
ingredients that the construction entangles. \emph{Polynomial activation} and
\emph{projective normalization} are the two axes of the study, and the
experiments support four conclusions.

\begin{enumerate}
\item \textbf{The polynomial nonlinearity alone is not competitive with the
modern defaults at equal parameter count.} At depth ten on MNIST, the bare
cubic $z + \lambda z^{3}$ lands almost four accuracy points below ReLU, GELU and
SiLU. The parameter-free projection closed a little over half of that gap, but
did not close it, and the remainder was not recovered by increasing the degree
or by making the polynomial learnable. Where the polynomial does matter is that
it responds strongly to scale control: the difference between the polynomial
\emph{with} and \emph{without} the projection is much larger than the difference
between any two polynomial orders or coefficients.

\item \textbf{Projective normalization, not the polynomial, drives the depth
behaviour.} The difference between the polynomial with and without the
projection is substantially larger than the difference between any two
polynomial orders or coefficients. The $(\ell_2)$ projection is what makes the
deep stack trainable: without it the polynomial dies at twenty layers, with it
the network survives to thirty, and at depth thirty it matches ReLU to within
$\LthirtyProjVsReluPt$ points --- with zero learned parameters and no batch
statistics. The benefit has a measured boundary: at depth fifty the projection
no longer suffices and every design in our study collapses.

\item \textbf{A parameter-free projection can stand in for BatchNorm and
LayerNorm at depth.} It matches their accuracy at depth while adding no
learned parameters and requiring no batch statistics, which makes it
attractive in settings where batch-statistics coupling is undesirable ---
small batches, streaming inference, or length-varying sequences.

\item \textbf{Gating is a parameter-budget question, and it requires
normalization.} Once parameter counts are matched, gated units confer no clear
accuracy advantage in these experiments. A gated unit fails
to train at depth ten without normalization, with activation norms decaying to
the point where the
output gradient vanishes, and normalization of any kind repairs it. Gated
units should therefore be read as components of a normalized block, not as
stand-alone improvements.

\end{enumerate}

The practical summary is simple, deliberately deflationary for the polynomial
and constructive for projective normalization: the cubic activation is not a
better nonlinearity than ReLU/GELU/SiLU, but controlling activation scale
across depth is first-order while the choice of activation is second-order. The
projective $\ell_2$ map is a cheap, parameter-free way to do the former, and it
can replace BatchNorm or LayerNorm in a deep stack without loss of accuracy.

Several directions follow naturally. Extending the comparison to convolutional
and attention architectures would test whether the projection's advantages
survive where residual connections and the attention softmax already control
scale. Repeating the normalization comparison under a per-configuration
learning-rate sweep would quantify how much of the observed ordering is
optimizer geometry rather than architecture. Finally, since learnable
polynomial activations (Legendre, Chebyshev, B-spline) recovered little beyond
the fixed cubic while costing appreciably more per step, a useful follow-up
would ask under what data regimes --- if any --- that extra flexibility is
worth its cost.

\paragraph{Reproducibility and code availability.}
Every number in this paper is produced by the harness in
\texttt{experiments/} and injected into the text as a \LaTeX{} macro by
\texttt{make\_macros.py}; the tables and figures are emitted by
\texttt{analyze.py}. The reference network is reproduced exactly, to the
printed precision of its own logs, by the harness configured with a depth-ten
projective polynomial and no separate normalization
(Table~\ref{tab:ref}).

The complete code, the per-configuration JSON logs behind every table and
figure, and the sources of this paper are released under the MIT license:

\begin{center}
\url{https://github.com/lfwin/Polynomial-Activation-and-projective-normalization}
\end{center}

\noindent
The repository contains the harness (\texttt{experiments/}), the analysis and
macro-generation scripts, the raw results, and the \LaTeX{} sources of this
document; the README describes how to regenerate every result from scratch.

\begin{table}[t]
\centering
\caption{Reproducing the reference network (depth-ten projective polynomial, no separate normalisation) with the unified harness. The published \texttt{results.json} reports $\refTestAcc\%$ test accuracy after $\refEpochs$ epochs. Epoch-one train/test values agree to the precision printed in the original logs; epoch two onward the trajectory depends on the annealing budget ($T_{\max}=\texttt{epochs}$), which differs between the 20-epoch reference run and the 2-epoch fidelity check.}
\label{tab:ref}
\small
\begin{tabular}{lrr}
\toprule
Run & Train (\%) & Test (\%) \\
\midrule
Reference, epoch 1 & $\refEoneTrain$ & $\refEoneTest$ \\
Reference, epoch 2 & $\refEtwoTrain$ & $\refEtwoTest$ \\
Reference, final & $\refTrainAcc$ & $\refTestAcc$ \\
\midrule
Harness (2\,ep), epoch 1 & $\fidEOneTrain$ & $\fidEOneTest$ \\
Harness (2\,ep), epoch 2 & $\fidETwoTrain$ & $\fidETwoTest$ \\
\bottomrule
\end{tabular}
\end{table}

\section*{Acknowledgements}

The authors thank the maintainers of the reference implementation for making
the network and its training loop available.
